\documentclass[11pt,a4paper]{article}

\usepackage[utf8]{inputenc}
\usepackage[T1]{fontenc}
\usepackage{amsmath,amssymb,amsthm,amsfonts}
\usepackage{mathtools}
\usepackage{graphicx}
\usepackage{booktabs}
\usepackage{hyperref}
\usepackage[margin=1in]{geometry}
\usepackage{bm}
\usepackage{physics}
\usepackage{xcolor}
\usepackage{caption}

\hypersetup{colorlinks=true, linkcolor=blue!60!black, citecolor=blue!60!black, urlcolor=blue!60!black}

\theoremstyle{definition}
\newtheorem{definition}{Definition}
\newtheorem{proposition}{Proposition}

\title{Differential Log-Negativity of the Measured \texorpdfstring{$t\bar{t}$}{ttbar} Spin Density Matrix, with a Topological Extension to Higgs Decays}

\author{Trevor Scott\thanks{Correspondence: \texttt{[your email]}. Code and data: \texttt{[github URL]}.}}
\date{\today}

\begin{document}
\maketitle

\begin{abstract}
We present the first computation of the logarithmic negativity $E_N(t{:}\bar{t})$ of the measured $t\bar{t}$ spin density matrix resolved across invariant-mass bins, using the full 15-coefficient reconstruction from CMS's 138\,fb$^{-1}$ differential measurement~\cite{CMS2024}. With covariance-matrix-propagated uncertainties and positivity-preserving projection of unphysical samples, we find $E_N(t{:}\bar{t}) = 0.18\,[0.13, 0.29]$ bits at the $m_{t\bar{t}}$ threshold (300--400\,GeV, 68\% CL), and $E_N$ consistent with zero in the three higher-mass bins (400--600, 600--800, $>$800\,GeV) at the resolution of the inclusive $|\cos\theta|$ selection. Alongside $E_N$ we compute both conventional singlet witness $D = -\text{Tr}(C)/3$ and the triplet witness $\tilde{D} = (C_{nn} - C_{kk} - C_{rr})/3$, observing the expected monotonic transition from singlet-like threshold region to triplet-like boosted regime. As a theoretical extension to multipartite systems, we apply the flag-complex persistent-homology construction of the Relational Witnessing Framework~\cite{Scott2026RWF} to the Standard Model tripartite density operator for $H \to ZZ^*$ of Aguilar-Saavedra~\cite{AS2024}, finding an $L$-centric topological signature with $E_N(L{:}S_i) \approx 1.0$ bits and $E_N(S_1{:}S_2) \lesssim 0.01$ bits. This extension is a prediction for Run 3, not a claim about existing measurement; the tripartite $H \to ZZ^*$ density operator has not yet been measured. All code and data are public; see Section~\ref{sec:code}.
\end{abstract}

\section{Introduction}

The observation of quantum entanglement in top-quark pairs by ATLAS~\cite{ATLAS2024} and CMS~\cite{CMS2024, CMS2024entangle} has established collider experiments as a new platform for quantum information measurements at TeV energy scales. These measurements go beyond the traditional kinematic program by reconstructing the reduced density matrix of the final-state quantum system from angular distributions and then computing entanglement measures directly from it.

The published entanglement claims use one of two scalar witnesses: the singlet witness $D = -\text{Tr}(C)/3$, where entanglement implies $D < -1/3$ for $t\bar{t}$ from unpolarized $pp$ production, or the triplet witness $\tilde{D} = (C_{nn} - C_{kk} - C_{rr})/3$, where entanglement implies $\tilde{D} > +1/3$. Both are inequalities that are sufficient for entanglement but not necessary; they are efficient experimentally because they collapse the full $4 \times 4$ density matrix to a single scalar.

In this paper we make two contributions. First, we compute the logarithmic negativity
\begin{equation}
E_N(\rho) = \log_2 \lVert \rho^{T_A} \rVert_1
\end{equation}
directly from the measured 15-coefficient $t\bar{t}$ spin density matrix, in each $m_{t\bar{t}}$ bin. This is a stronger observable than $D$ or $\tilde{D}$: it is zero exactly when $\rho$ is separable (by the Peres-Horodecki criterion), it is additive, and it quantifies the entanglement on an information-theoretic scale (bits). We propagate uncertainties via Monte Carlo on the full 16-dimensional covariance matrix from HEPData, with positivity-preserving projection of unphysical samples. Second, we extend the methodology to multipartite quantum systems by applying the flag-complex persistent-homology construction of~\cite{Scott2026RWF} to the Standard Model tripartite density operator of~\cite{AS2024} for $H \to ZZ^*$. This section is a theoretical prediction for a system that has not yet been measured, setting a target for Run 3 tripartite tomography.

\section{Methodology}\label{sec:methodology}

\subsection{Density matrix reconstruction}

For $t\bar{t}$ in the helicity basis $\{\hat{k}, \hat{n}, \hat{r}\}$, the density matrix has the Fano-Bloch decomposition
\begin{equation}
\rho_{t\bar{t}} = \frac{1}{4}\!\left(I \otimes I + \vec{B}_+ \cdot \vec{\sigma} \otimes I + I \otimes \vec{B}_- \cdot \vec{\sigma} - \sum_{ij} C_{ij}\, \sigma_i \otimes \sigma_j\right). \label{eq:fano}
\end{equation}
CMS publishes the 6 polarization components $\vec{B}_\pm$ and the 9 elements of $C$, decomposed into a symmetric $C_{ij}^+ = (C_{ij} + C_{ji})/2$ and antisymmetric $C_{ij}^- = (C_{ij} - C_{ji})/2$ set. We use the ATLAS/CMS sign convention (negative $C$ in the Fano decomposition) throughout, and all 15 (+ normalization) coefficients simultaneously. The negative sign on $C$ means that same-helicity top-antitop pairs have positive $C_{kk}$ and that the spin-singlet state has $C = +I$.

\subsection{Log-negativity and entanglement witnesses}

The logarithmic negativity is computed by taking the partial transpose $\rho^{T_t}$ on the top subsystem, computing its trace norm as the sum of absolute values of the eigenvalues, and taking $\log_2$. For comparison we also compute
\begin{align}
D &= -(C_{kk} + C_{nn} + C_{rr})/3, & \text{entangled if } D < -\tfrac{1}{3}, \\
\tilde{D} &= (C_{nn} - C_{kk} - C_{rr})/3, & \text{entangled if } \tilde{D} > +\tfrac{1}{3}.
\end{align}
$D$ witnesses singlet-like states (all $C_{ii}$ positive, $\text{Tr}(C)$ large positive). $\tilde{D}$ witnesses triplet-like states at boosted kinematics, where parity and rotation structure cause $C_{kk}, C_{rr}$ to flip sign relative to $C_{nn}$. Neither witness is sharper than $E_N$; but they are the scalar observables CMS reports directly.

\subsection{Uncertainty propagation}\label{sec:uncertainty}

The CMS differential measurement publishes both central values with stat+syst uncertainties (HEPData table t9 of record \texttt{ins2829523}) and the full 64$\times$64 covariance matrix across the four $m_{t\bar{t}}$ bins (table t10). We restrict to the intra-bin 16$\times$16 blocks and sample coefficient vectors from the multivariate Gaussian with the measured mean and covariance.

For each sample we reconstruct $\rho$ via~\eqref{eq:fano} and compute the partial-transpose eigenvalues. When the measured central values or a drawn sample yield a $\rho$ with negative eigenvalues (a known feature of unconstrained coefficient tomography, since the 15 measurements are independent and not required to produce a positive semidefinite matrix), we project onto the positive-semidefinite cone by clipping the eigenvalues at zero and renormalizing. This preserves the Hermiticity and unit trace of $\rho$ and, crucially, never inflates $E_N$: clipping only removes unphysical mass that would otherwise contribute to (or from) negative PT eigenvalues. We report the 16--84 percentile of $E_N$ across 20{,}000 samples as the 68\% CL interval.

Positivity violation rates in our MC are 40\% in the threshold bin, 54\% in the 400--600 bin, 6\% in the 600--800 bin, and 0\% in the boosted bin, reflecting the increasing statistical precision at higher mass. We verify that the quoted $E_N$ intervals are robust to the choice of projection method.

\section{Measured $t\bar{t}$ results}\label{sec:ttbar}

\begin{figure}[t]
\centering
\includegraphics[width=\textwidth]{fig_ttbar_3panel.png}
\caption{Quantum entanglement observables for the measured $t\bar{t}$ state across $m(t\bar{t})$ bins, computed from the CMS 138\,fb$^{-1}$ differential measurement~\cite{CMS2024}. \textbf{Left:} log-negativity $E_N$ reconstructed from the full 15-coefficient spin density matrix with covariance-propagated 68\% CL error bars. \textbf{Center:} singlet witness $D$ with the entanglement threshold $D = -1/3$ (green dashed). \textbf{Right:} triplet witness $\tilde{D}$ with its threshold $\tilde{D} = +1/3$ (purple dashed). The state transitions from singlet-like entangled at threshold ($D < -1/3$, $E_N > 0$) to triplet-like non-entangled (in this inclusive $|\cos\theta|$ selection) at high mass, with $\tilde{D}$ rising monotonically toward the triplet entanglement boundary.}
\label{fig:3panel}
\end{figure}

Table~\ref{tab:results} summarizes the per-bin results; Figure~\ref{fig:3panel} shows the three-panel visualization.

\begin{table}[t]
\centering
\begin{tabular}{lrrr}
\toprule
$m_{t\bar{t}}$ range [GeV] & $D$ & $\tilde{D}$ & $E_N$ [bits, 68\% CL] \\
\midrule
$300 < m_{t\bar{t}} < 400$   & $-0.412 \pm 0.035$ & $-0.052 \pm 0.036$ & $\bm{0.18\,[0.13, 0.29]}$ \\
$400 < m_{t\bar{t}} < 600$   & $-0.204 \pm 0.012$ & $-0.013 \pm 0.012$ & $0$ \\
$600 < m_{t\bar{t}} < 800$   & $-0.089 \pm 0.020$ & $+0.069 \pm 0.020$ & $0$ \\
$m_{t\bar{t}} > 800$         & $-0.005 \pm 0.024$ & $+0.112 \pm 0.025$ & $0$ \\
\bottomrule
\end{tabular}
\caption{CMS-measured observables per $m_{t\bar{t}}$ bin. Entanglement is detected by $E_N > 0$ and by $D < -1/3$ in the threshold bin. $\tilde{D}$ rises monotonically toward (but does not cross) the triplet entanglement threshold $+1/3$ in this inclusive selection; CMS observes boosted entanglement at $>5\sigma$ with the additional $|\cos\theta| < 0.4$ cut~\cite{CMS2024}. Uncertainty intervals on $E_N$ are from 20{,}000-sample MC on the full 16$\times$16 intra-bin covariance; intervals on $D$ and $\tilde{D}$ combine stat and syst quadratically.}
\label{tab:results}
\end{table}

\subsection{Discussion}

\textbf{Threshold entanglement.} In the $300$--$400$ GeV bin, the near-singlet state yields $E_N = 0.18\,[0.13, 0.29]$ bits at 68\% CL, corresponding to a most-negative PT eigenvalue of $-0.067$. This is a genuine information-theoretic measure of the entanglement, stronger than the $D < -1/3$ observation. The uncertainty budget is dominated by the statistics on the diagonal $C$ coefficients, which have stat errors of order 6--7\%.

\textbf{Singlet--triplet transition.} Across the four $m_{t\bar{t}}$ bins, $D$ rises monotonically from $-0.412$ at threshold to $-0.005$ in the boosted bin, while $\tilde{D}$ rises from $-0.052$ to $+0.112$. The crossover of $\tilde{D}$ through zero occurs between 400 and 600 GeV and reflects the sign change in $C_{kk}$ and $C_{rr}$ that is characteristic of the transition from s-channel-dominated singlet production at threshold to the boosted kinematic regime where parity structure dominates. This is the first time this evolution has been presented with log-negativity as an observable alongside the scalar witnesses.

\textbf{Absence of $E_N$ at high mass.} In the bins above 400 GeV, our reconstructed $E_N$ is consistent with zero at the full resolution of the measurement. CMS's own boosted entanglement observation at $>5\sigma$~\cite{CMS2024} applies an additional $|\cos\theta| < 0.4$ selection that we have not imposed here; that selection enhances the triplet-like contribution and lifts $\tilde{D}$ above $+1/3$. Extending this analysis to the $|\cos\theta|$-selected subsample (HEPData table t11 of the same record) is a straightforward next step.

\textbf{Comparison to Standard Model.} Our measured central values are consistent with the POWHEG+Pythia8, POWHEG+Herwig7, MadGraph5\_aMC@NLO+Pythia8, and MiNNLO+Pythia8 predictions published alongside the CMS data, within uncertainties, across all four bins.

\section{Theoretical extension: $H \to ZZ^*$ tripartite topology}\label{sec:hzz}

We now apply the flag-complex construction of the Relational Witnessing Framework~\cite{Scott2026RWF} to the tripartite quantum state produced in $H \to ZZ^*$ decays. This section is a theoretical prediction for a system that has not been measured: Aguilar-Saavedra~\cite{AS2024} recently showed that the Higgs-to-diboson decay produces a genuinely tripartite state in the Hilbert space $\mathcal{H}_L \otimes \mathcal{H}_{S_1} \otimes \mathcal{H}_{S_2}$ (dimensions $9 \times 3 \times 3$) and that Run 3 data will measure its density operator at $>5\sigma$ significance.

\subsection{The flag-complex construction}

Given an $n$-partite state $\rho$, RWF Definition 2 constructs the entanglement complex $K_W$ as:
\begin{itemize}
\item 0-simplices: the $n$ subsystems;
\item 1-simplices: pairs $(i, j)$ with weight $w_{ij} = E_N(i{:}j) = \log_2 \lVert \rho_{ij}^{T_i}\rVert_1$ on the reduced pair state $\rho_{ij} = \text{Tr}_{\{k:k\neq i,j\}}(\rho)$;
\item $k$-simplices: subsets $\{i_1, \ldots, i_{k+1}\}$ are included iff all pairwise edges are present (flag closure).
\end{itemize}
Persistent homology on the decreasing-$\epsilon$ filtration yields the topological signature $\chi(K_W) = \{\text{Dgm}_k(K_W)\}_k$.

For the bipartite $t\bar{t}$ case of Section~\ref{sec:ttbar}, $K_W$ has one edge; the construction reduces to the scalar $E_N$ reported in Table~\ref{tab:results}. The framework's content appears at $n \geq 3$.

\subsection{$H \to ZZ^*$ Standard Model prediction}

We construct $|\psi\rangle_{L S_1 S_2}$ using Aguilar-Saavedra's helicity-amplitude decomposition~\cite{AS2024}, parameterized by three complex amplitudes $a_{11}, a_{00}, a_{-1-1}$ fixed by the SM at leading order by
\begin{equation}
a_{11} = a_{-1-1} = M_Z \cdot m_{Z^*}, \quad a_{00} = \tfrac{1}{2}(M_H^2 - M_Z^2 - m_{Z^*}^2),
\end{equation}
with $m_{Z^*}$ the invariant mass of the off-shell $Z$. We scan $m_{Z^*} \in [5, 70]$ GeV and compute the pairwise log-negativities.

\begin{figure}[t]
\centering
\includegraphics[width=0.85\textwidth]{fig1_hzz_edge_weights.png}
\caption{SM prediction for the pairwise log-negativities of the $H \to ZZ^*$ tripartite state as a function of $m_{Z^*}$. The two $L$-$S$ edges are identical by symmetry and near $1$ bit; the $S_1$-$S_2$ edge is two orders of magnitude smaller.}
\label{fig:hzz_edges}
\end{figure}

We find a distinctive \emph{$L$-centric} topological signature:
\begin{itemize}
\item $E_N(L{:}S_1) = E_N(L{:}S_2) \approx 1.0$--$1.07$ bits across the range, decreasing slowly with $m_{Z^*}$;
\item $E_N(S_1{:}S_2) \leq 0.011$ bits, two orders of magnitude smaller;
\item The 2-simplex $\{L, S_1, S_2\}$ enters the flag complex only at $\epsilon \lesssim 0.01$; for most of the filtration, the 1-skeleton is $V$-shaped.
\end{itemize}
This is a direct topological encoding of the fact that in $H \to ZZ^*$, the two $Z$-spins are primarily entangled through the orbital angular momentum channel rather than with each other directly: tracing out $L$ from the tripartite state leaves a near-separable $S_1$-$S_2$ reduced density matrix.

\subsection{Limitations}

This section is a theoretical prediction. The Aguilar-Saavedra density operator is an SM leading-order analytical construction; the tripartite operator has not yet been measured. With Run 3 $H \to ZZ \to 4\ell$ data, the tripartite tomography of~\cite{AS2024} establishes the measured operator at $>5\sigma$ per bin, at which point this analysis could be repeated on measured data and reported alongside the bipartite $t\bar{t}$ result as a genuinely multipartite collider observable.

We note two limitations of the flag-complex construction itself, also discussed in~\cite{Scott2026RWF}: (i) the flag complex is blind to genuinely multipartite entanglement that leaves no trace on pair marginals. We have verified numerically that $n$-qubit GHZ states and $4$-qubit cluster states both yield trivial flag complexes. The $H \to ZZ^*$ state, by contrast, produces the non-trivial $L$-centric signature we report because its multipartite structure does leak through pair marginals. (ii) The $l=0$ Clebsch-Gordan block in our construction of $|\psi\rangle$ uses a $1/\sqrt{3}$ ansatz that we have not verified against~\cite{AS2024}'s eq.~(7). The $l=2$ block, which dominates $E_N(L{:}S_i)$, is exact; the residual $S_1$-$S_2$ entanglement may shift by an $\mathcal{O}(1)$ factor but remains suppressed relative to the $L$-$S$ edges.

\section{Code availability}\label{sec:code}

All code and data are publicly available:

\begin{itemize}
\item \textbf{Pipeline}: density matrix reconstruction, log-negativity computation, flag-complex persistent homology, covariance-propagated Monte Carlo.
\item \textbf{CMS data}: from HEPData record \texttt{ins2829523}~\cite{CMS2024data} (tables t9, t10).
\item \textbf{Repository}: \texttt{[github URL]}.
\end{itemize}

The analysis reproduces from a clean Python environment in under 30 seconds per script.

\section{Conclusion}

We have presented the first log-negativity measurement of the $t\bar{t}$ spin density matrix at the LHC, resolved across the $m_{t\bar{t}}$ phase space, with covariance-propagated uncertainties. The measured threshold state has $E_N = 0.18\,[0.13, 0.29]$ bits at 68\% CL. The complementary evolution of $D$ and $\tilde{D}$ across the four kinematic bins captures the singlet-to-triplet transition characteristic of QCD top-pair production. The framework extends to multipartite systems via the flag-complex construction of~\cite{Scott2026RWF}, and we report the SM prediction for $H \to ZZ^*$ as a target for Run 3 tripartite tomography.

\paragraph{Acknowledgments.} T.S. thanks [collaborators] for discussion on quantum observables at the LHC and on the Aguilar-Saavedra construction.

\begin{thebibliography}{99}

\bibitem{Scott2026RWF}
T.~Scott, ``The Relational Witnessing Framework: The Emergence of Time and Experience from Entanglement Entropy,'' (2026).

\bibitem{AS2024}
J.~A.~Aguilar-Saavedra, ``Tripartite entanglement in $H \to ZZ, WW$ decays,'' \emph{Phys.~Rev.~D} \textbf{109}, 113004 (2024) [\texttt{arXiv:2403.13942}].

\bibitem{ATLAS2024}
ATLAS Collaboration, ``Observation of quantum entanglement with top quarks at the ATLAS detector,'' \emph{Nature} \textbf{633}, 542 (2024) [\texttt{arXiv:2311.07288}].

\bibitem{CMS2024}
CMS Collaboration, ``Measurements of polarization and spin correlation and observation of entanglement in top quark pairs using lepton+jets events from proton-proton collisions at $\sqrt{s} = 13$ TeV,'' \emph{Phys.~Rev.~D} \textbf{110}, 112016 (2024) [\texttt{arXiv:2409.11067}].

\bibitem{CMS2024data}
CMS Collaboration, HEPData record \texttt{ins2829523} (2024), \url{https://doi.org/10.17182/hepdata.153301.v1}.

\bibitem{CMS2024entangle}
CMS Collaboration, ``Observation of quantum entanglement in top quark pair production in proton--proton collisions at $\sqrt{s} = 13$ TeV,'' \emph{Rep.~Prog.~Phys.} \textbf{87}, 117801 (2024) [\texttt{arXiv:2406.03976}].

\bibitem{Afik2021}
Y.~Afik and J.~R.~M.~de~Nova, ``Entanglement and quantum tomography with top quarks at the LHC,'' \emph{Eur.~Phys.~J.~Plus} \textbf{136}, 907 (2021) [\texttt{arXiv:2003.02280}].

\bibitem{Peres1996}
A.~Peres, ``Separability criterion for density matrices,'' \emph{Phys.~Rev.~Lett.} \textbf{77}, 1413 (1996).

\bibitem{Vidal2002}
G.~Vidal and R.~F.~Werner, ``Computable measure of entanglement,'' \emph{Phys.~Rev.~A} \textbf{65}, 032314 (2002).

\bibitem{Gudhi}
The GUDHI Project, \emph{GUDHI User and Reference Manual} (GUDHI Editorial Board, 2024), \url{https://gudhi.inria.fr/}.

\bibitem{CohenSteiner2007}
D.~Cohen-Steiner, H.~Edelsbrunner, and J.~Harer, ``Stability of persistence diagrams,'' \emph{Discrete Comput.~Geom.} \textbf{37}, 103 (2007).

\end{thebibliography}

\end{document}
